\documentclass[11pt,a4paper]{article}

\usepackage[utf8]{inputenc}
\usepackage[T1]{fontenc}
\usepackage{lmodern}
\usepackage{amsmath,amssymb}
\usepackage{booktabs}
\usepackage{array}
\usepackage{graphicx}
\usepackage{float}
\usepackage[margin=1in]{geometry}
\usepackage{hyperref}
\graphicspath{{figures/}}

\sloppy
\emergencystretch=3em

\newcommand{\avert}{a_{\mathrm{vert}}}
\newcommand{\amag}{\lvert a\rvert}
\newcommand{\ghat}{\hat{g}}
\newcommand{\ahat}{\hat{a}}
\newcommand{\zhat}{\hat{z}}

\title{Accelerometer Algorithm Inputs --- What Signal Each Stage Matches Against}
\author{}
\date{\today}

\begin{document}
\maketitle

\paragraph{Purpose.}
Quick reference for what every accelerometer-only algorithm in this
repo consumes as input, what it does in plain words, and how the
trapezoid match in \emph{segmentation} differs from the trapezoid
match in \emph{prediction}. Barometer paths are out of scope.

% ===========================================================
\section{Common signal definitions}
% ===========================================================

Every accelerometer algorithm starts from the raw 3-axis device-frame
sample $(a_x, a_y, a_z)$ at sampling rate $f_s$ (typically 100\,Hz).
From it we derive three candidate scalar signals:

\begin{itemize}\setlength\itemsep{2pt}
  \item $\amag \;=\; \sqrt{a_x^2 + a_y^2 + a_z^2}$ \\
        Gravity \emph{included}. Orientation-blind.\\
        \textbf{Used by:} the GT editor (\texttt{src/data/gt\_editor.py:227}) for visual annotation only. \textbf{No active algorithm consumes $\amag$.}
  \item $a_z$ \\
        Device-frame $z$-axis. \textbf{No active algorithm consumes $a_z$.}
  \item $\avert \;=\; (a - \hat g) \cdot \dfrac{\hat g}{\lVert \hat g\rVert}$ \\
        Signed scalar projection of the acceleration vector onto the
        estimated gravity direction $\hat g$. Positive = up. \\
        \textbf{This is the signal every active accelerometer algorithm
        uses.}
\end{itemize}

The gravity vector $\hat g$ is estimated from a stationary window:
\texttt{estimate\_gravity\_stationary} in
\texttt{src/utils/accelerometer\_utils.py}. In segmentation the
window is the first 0.5\,s of the buffer
(\texttt{...template\_match/fit\_elevator\_parameters/common.py:136}).
In prediction the pre/post ride stationary windows are preferred,
falling back to in-ride estimation, then to $\amag$ as a last resort
(\texttt{trapezoid\_accel/estimator.py:118--170}).

After projection the signal is detrended (slow rolling mean) and
smoothed (rolling mean over \texttt{smooth\_sec}=0.4\,s) to remove
DC bias from orientation drift and hand-tremor:
\[
\avert^{\text{(smooth)}}(t) \;=\; \mathrm{RollMean}\bigl(\avert(t) - \mathrm{RollMean}_{\text{slow}}(\avert(t)),\; w_{\text{smooth}}\bigr).
\]

% ===========================================================
\section{Segmentation --- accelerometer path}
% ===========================================================

Single active algorithm: \textbf{trapezoid matched-filter template match}, dispatched as
\verb|SegmentAlgorithm.ACC_TEMPLATE_MATCH|.\\
Files:
\texttt{src/segmentation/algorithms/accelerometer\_only/template\_match/}\,---\,
\texttt{matcher.py} (legacy learned-template path) and
\texttt{check\_grid\_across\_signal/detect.py} (active synthetic-grid path,
wired to the editor UI).

\paragraph{Input.}
Smoothed $\avert$ over the \emph{entire session}
(\texttt{matcher.py:118}, \texttt{detect.py} via the shared
preprocessing in
\texttt{fit\_elevator\_parameters/common.py:135--144}). The matcher
also computes an LPF velocity $v_{\text{lpf}}$ by integrating
$\avert$, and the two scores are linearly fused
(\texttt{matcher.py:125--127}); the active grid detector uses
$\avert$ alone.

\paragraph{Template.}
Synthetic trapezoid kernel parameterised by half-width $W$ and
flat-fraction $f \in [0,1]$:
\[
\text{tpl}(t;\,t_c,W,f) \;=\;
\begin{cases}
1 & |t-t_c| \le f\,W \quad\text{(flat top)} \\[2pt]
\dfrac{W - |t-t_c|}{W\,(1-f)} & f\,W < |t-t_c| \le W \quad\text{(linear ramp)}\\[6pt]
0 & |t-t_c| > W.
\end{cases}
\]
Grid: $30 \times 15 = 450$ templates over $W \in [0.3,\,3.0]$\,s,
$f \in [0.05,\,0.80]$, plus the $f=0$ triangle row
(\texttt{detect.py:130--143}).

\paragraph{What it does (problem type: \emph{detection}).}
Slide every template across the whole signal using Pearson normalised
cross-correlation (\texttt{ncc\_slide}, \texttt{matcher.py:33--54}).
Reduce the per-sample argmax across $(W,f)$ to a peak score, then:
\begin{enumerate}\setlength\itemsep{2pt}
  \item \textbf{Peak-pick} on unsigned $R^2$ with threshold
        \texttt{r2\_peak\_thresh}=0.40 and a minimum matched
        amplitude $|A|\ge 0.20\,\mathrm{m/s^2}$.
  \item \textbf{NMS dedup} within \texttt{nms\_radius\_s}=1.0\,s and
        same-sign suppression within 5\,s.
  \item \textbf{Pair} a positive (take-off) peak with a later
        negative (landing) peak, accept if their shared-shape
        joint $R^2 \ge 0.90$, the pair $|A|\ge 0.30\,\mathrm{m/s^2}$,
        and the plateau between the lobes is quiet
        (\texttt{quiet\_middle\_ratio}=0.5).
\end{enumerate}

\paragraph{Output.}
A list of ride intervals $(t_{\text{start}}, t_{\text{end}}, \text{direction})$. No height value is produced.

% ===========================================================
\section{Prediction --- accelerometer paths}
% ===========================================================

Two accelerometer-only predictors share the
\texttt{predict\_segment(ride, pre, post) -> PredictionOutput}
contract dispatched by \texttt{Predictor} in
\texttt{src/prediction/algorithms/predictor.py:29--43}.

% -------------------------------------------------
\subsection{ZUPT double-integration}
% -------------------------------------------------
File: \texttt{src/prediction/algorithms/accelerometer\_only/zupt\_accel/estimator.py}.

\paragraph{Input.}
$\avert$ over the \emph{ride window}, plus the pre/post stationary
windows used to estimate $\hat g$.

\paragraph{What it does.}
\begin{enumerate}\setlength\itemsep{2pt}
  \item Detect the active-motion sub-window by thresholding
        $|\avert|$ (\texttt{\_find\_motion\_window}, \texttt{estimator.py:50--65}).
  \item Zero-force $\avert(t)=0$ outside that sub-window
        (\texttt{\_zupt\_integrate\_windowed}, \texttt{estimator.py:68--79}).
  \item Double-integrate $\avert \to v \to h$, applying linear-drift
        correction so $v$ returns to zero at the boundaries.
\end{enumerate}

\paragraph{Trapezoid involvement.}
\textbf{None.} ZUPT does not match any template; it is a pure
double-integration with zero-velocity-update bookends.

\paragraph{Output.}
Signed $\Delta h$ in metres plus a theoretical $\sigma$ from white-noise integration bounds.

% -------------------------------------------------
\subsection{Trapezoid pulse-pair fit}
% -------------------------------------------------
File: \texttt{src/prediction/algorithms/accelerometer\_only/trapezoid\_accel/estimator.py}; kernel + grid search in \texttt{pulse\_pair.py}.

\paragraph{Input.}
Detrended + smoothed $\avert$ over the \emph{ride window}
(\texttt{estimator.py:190--196}): same projection, same
\texttt{smooth\_sec}=0.4\,s rolling mean.

\paragraph{What it does (problem type: \emph{parameter estimation}).}
Searches the same $(W,f)$ trapezoid grid as segmentation
(\texttt{pulse\_pair.py:42}, $30\times 15$). For every grid point and
every candidate centre pair $(t_{c1}, t_{c2})$ it solves
least-squares for a \emph{shared} amplitude $|A|$ across the two
lobes and scores by the joint $R^2 = \tfrac12(R^2_1 + R^2_2)$
(\texttt{pulse\_pair.py:42--62}). The argmax across
$(W, f, t_{c1}, t_{c2})$ is kept. A second 4-parameter
\emph{joined-pulse} fit ($\Delta t_c = 2W$ forced) is run in parallel;
the better joint-$R^2$ wins (\texttt{estimator.py:213--254}). A ZUPT
fallback kicks in if both matched-filter fits collapse.

\paragraph{Output.}
Closed-form
\[
\Delta h \;=\; \operatorname{sign}\cdot |A|\cdot W\cdot (1+f)\cdot (t_{c2} - t_{c1})
\]
(\texttt{estimator.py:330--339}), plus a delta-method theoretical
$\sigma$ and a conformal-calibrated 90\,\% CI half-width.

% ===========================================================
\section{Segmentation vs.\ prediction --- same trapezoid, different question}
% ===========================================================

Both stages match the same trapezoid template family against the same
projected signal $\avert$. The difference is what they do with the
match.

\begin{table}[h]
\centering
\renewcommand{\arraystretch}{1.15}
\begin{tabular}{p{0.22\linewidth}p{0.36\linewidth}p{0.36\linewidth}}
\toprule
 & \textbf{Segmentation} & \textbf{Prediction (trapezoid pulse-pair)} \\
\midrule
\textbf{Input scope} & smoothed $\avert$ over the \textbf{whole session} & detrended+smoothed $\avert$ over \textbf{one ride window} \\
\textbf{Template family} & synthetic trapezoid $(W,f)$ grid $+$ $f{=}0$ triangle row & same $(W,f)$ grid, shared-shape constraint across the lobe pair \\
\textbf{Scoring} & sliding Pearson NCC, per-sample peak score & exhaustive LS amplitude per $(W,f,t_{c1},t_{c2})$, joint $R^2$ \\
\textbf{What is found} & peak \emph{locations} (when does a ride happen?) & best-fit \emph{parameters} $(W,f,|A|,t_{c1},t_{c2})$ \\
\textbf{Problem class} & \textbf{detection} & \textbf{parameter estimation} \\
\textbf{Output} & ride intervals $(t_s, t_e, \text{dir})$ & $\Delta h$ in metres $+$ CI half-width \\
\textbf{Filters / gates} & peak $R^2 \ge 0.40$, NMS, pair $R^2 \ge 0.90$, quiet-middle & quality score, $R^2$ thresholds by ride length, anchor-ratio, $\Delta t_c$ floors \\
\bottomrule
\end{tabular}
\end{table}

The trapezoid kernel itself lives in a single shared utility
(\texttt{src/utils/trapezoid\_template.py}, re-exported by both
stages) so the template parametrisation is identical by construction.

% ===========================================================
\section{One-line summary}
% ===========================================================

\textit{All active accelerometer algorithms --- the segmentation
matched-filter, the ZUPT predictor, and the trapezoid pulse-pair
predictor --- consume $\avert$ (gravity-projected vertical
acceleration), not $\amag$.} The trapezoid template appears in both
stages, but with different roles: segmentation uses it as a
\emph{detector} (sliding NCC over the whole session to find rides),
and prediction uses it as an \emph{estimator} (exhaustive LS fit
inside a single ride to recover $\Delta h$).

\paragraph{Implication for editor graphs.}
The GT editor (\texttt{src/data/gt\_editor.py}) currently plots
$\amag$, which no active algorithm uses. The segmentation editor
(\texttt{src/segmentation/algorithms/editor.py:118--121}) and the
Streamlit segmentation step
(\texttt{src/pipelines/streamlit/step3\_segmentation.py:199--209})
both already plot $\avert$ (with its smoothed overlay), which is the
signal the matched filter actually scores against. Harmonising the
three to display $\avert$ is a follow-up; this report is the
reference for the decision.

% ===========================================================
\section{Vertical-acceleration reconstruction from accel + gyro}
% ===========================================================

\paragraph{Motivation.}
Every algorithm above consumes $\avert = (a-\ghat)\cdot\hat\ghat$, the
accelerometer projected onto a \emph{single} gravity vector $\ghat$
estimated from the pre/post stationary windows
(\texttt{estimate\_gravity\_stationary}). That assumes the phone's
orientation is fixed for the whole ride. When the phone rotates mid-ride,
$\ghat$ is stale and a slice of the $9.81\,\mathrm{m/s^2}$ gravity term
leaks into $\avert$ as a time-varying artefact that DC/drift removal cannot
kill --- it blurs the trapezoid and, because ZUPT double-integrates, a
$5^\circ$ tilt error ($\approx 0.04\,\mathrm{m/s^2}$) becomes
$\approx 1.9\,\mathrm{m}$ over a $10\,\mathrm{s}$ ride. Reconstruction
tracks the gravity direction $\ghat(t)$ with the gyroscope instead.

\paragraph{Pipeline.}
Code: \texttt{src/physics/reconstruct\_az/}, invoked as the first line of
\texttt{Predictor.predict} / \texttt{Segmenter.detect} via the one-liner
\verb|reconstruct_az(type, acc, gyro)| (config field \texttt{reconstruct};
\texttt{"none"} = today's behaviour). Each filter estimates a body-to-world
orientation quaternion $q(t)$, seeded from the opening stationary window,
then
\[
  a_{\mathrm{world}}(t) = R\!\big(q(t)\big)\,a(t),
  \qquad
  \avert(t) = a_{\mathrm{world},z}(t) - g .
\]
To keep the downstream gravity projection unchanged, the predictors/detector
are actually fed the \emph{virtually-flat-phone} signal
$\big(a_{\mathrm{world},x},\,a_{\mathrm{world},y},\,a_{\mathrm{world},z}+g\big)$
--- gravity kept on $+z$ --- so \texttt{estimate\_gravity\_stationary} still
recovers $\lvert g\rvert\!\approx\!9.81$ and the projection returns the true
signed vertical accel.

\paragraph{The five filters.}
All run in IMU mode (accel $a$, gyro $\omega$, both body-frame; no
magnetometer --- yaw is irrelevant for the vertical channel). Shared
quantities: normalised accel $\ahat=a/\amag$; predicted gravity direction in
the body frame $\hat v = R(q)^\top\zhat$; and an \emph{accel-trust} weight
$w=\exp\!\big(-\tfrac12((\amag-g)/(\tau g))^2\big)$ that down-weights the
accelerometer correction during the ride's acceleration pulses (when
$\amag\neq g$), so ride acceleration is never mistaken for tilt.
\begin{enumerate}\setlength\itemsep{3pt}
  \item \textbf{Complementary.} Gyro integration
        $q\!\leftarrow\! q\otimes\exp(\tfrac12\omega\,\Delta t)$, then a
        partial rotation of $\ahat$ onto $\zhat$:
        $q\leftarrow \mathrm{slerp}\big(1,\delta q;\,\alpha w\big)\otimes q$,
        with $\delta q$ the shortest arc from $R(q)\ahat$ to $\zhat$.
  \item \textbf{Mahony} (PI on $SO(3)$, gyro-bias). Error
        $e=\ahat\times\hat v$; bias $b\leftarrow b-k_I\,w\,e\,\Delta t$;
        $\dot q=\tfrac12 q\otimes(\omega-b+k_P\,w\,e)$.
  \item \textbf{Madgwick} (gradient descent). With
        $f=R(q)^\top\zhat-\ahat$ and Jacobian $J$,
        $\dot q=\tfrac12 q\otimes\omega-\beta w\,\dfrac{J^\top f}{\lVert J^\top f\rVert}$.
  \item \textbf{Valenti / AQUA} (algebraic + adaptive gain). Predict with
        the gyro, form the algebraic $\Delta q_{\mathrm{acc}}$ (shortest arc
        $R(q)\ahat\!\to\!\zhat$), and apply it with a gain that ramps
        \emph{linearly to zero} once $\lvert\amag-g\rvert/g$ exceeds a
        threshold: $q\leftarrow\mathrm{slerp}(1,\Delta q_{\mathrm{acc}};\alpha)\otimes q$.
  \item \textbf{ESKF} (error-state Kalman). Nominal $(q,b)$ propagated by the
        gyro; error state $(\delta\theta,\delta b)$ tracked with covariance
        $P$. Accelerometer update uses $H=[\,[\hat v]_\times\;\; 0\,]$,
        measurement noise $R=(\sigma_a/w)^2 I$, then injects
        $q\leftarrow q\otimes[1,\tfrac12\delta\theta]$ and resets. Estimating
        $b$ online holds attitude through long low-excitation cruise phases.
\end{enumerate}
Primary references (\texttt{src/physics/reconstruct\_az/references/}):
Euston et al.\ 2008 (complementary), Mahony et al.\ 2008, Madgwick 2010,
Valenti et al.\ 2015, Sol\`a 2017.

\paragraph{Does it help? (measured).}
Tables~\ref{tab:recon-pred} and~\ref{tab:recon-seg} score every variant on
the real elevator set, produced by
\texttt{scripts/reconstruction\_comparison\_report.py}. Reconstruction
sharply reduces prediction error (best filter in \textbf{bold}) and removes
segmentation false positives; the gain is largest on rides where the phone
was rotated (Fig.~\ref{fig:recon-diag}).

\begin{table}[H]
\centering
\caption{Prediction accuracy by reconstruction (ZUPT and trapezoid).}
\label{tab:recon-pred}
\begin{tabular}{lrrrrr}
\toprule
\multicolumn{6}{l}{\textbf{ZUPT}} \\
\midrule
Reconstruction & MAE (m) & RMSE (m) & Cov@90 & $\le$1.5\,m & med CI (m) \\
\midrule
none & 2.545 & 18.827 & 0.919 & 0.758 & 2.256 \\
Complementary & 1.075 & 3.246 & 0.911 & 0.871 & 1.252 \\
\textbf{Mahony} & \textbf{1.068} & \textbf{3.238} & \textbf{0.916} & \textbf{0.869} & \textbf{1.237} \\
Madgwick & 1.476 & 10.096 & 0.924 & 0.868 & 1.258 \\
Valenti & 1.075 & 3.243 & 0.913 & 0.869 & 1.215 \\
ESKF & 1.083 & 3.074 & 0.908 & 0.874 & 1.165 \\
\bottomrule
\end{tabular}
\vspace{1em}

\begin{tabular}{lrrrrr}
\toprule
\multicolumn{6}{l}{\textbf{TRAPEZOID}} \\
\midrule
Reconstruction & MAE (m) & RMSE (m) & Cov@90 & $\le$1.5\,m & med CI (m) \\
\midrule
none & 1.734 & 11.313 & 0.835 & 0.848 & 0.894 \\
Complementary & 0.988 & 2.853 & 0.863 & 0.879 & 0.892 \\
\textbf{Mahony} & \textbf{0.943} & \textbf{2.734} & \textbf{0.864} & \textbf{0.880} & \textbf{0.892} \\
Madgwick & 1.173 & 3.452 & 0.850 & 0.871 & 0.892 \\
Valenti & 0.959 & 2.744 & 0.861 & 0.879 & 0.892 \\
ESKF & 0.964 & 2.786 & 0.866 & 0.884 & 0.892 \\
\bottomrule
\end{tabular}
\vspace{1em}

\end{table}

\begin{table}[H]
\centering
\caption{Segmentation accuracy by reconstruction (ACC template match).}
\label{tab:recon-seg}
\begin{tabular}{lrrrrrrrr}
\toprule
Reconstruction & $n_{gt}$ & detect & miss & FP & F1* & IoU-F1 & recall & prec \\
\midrule
none & 627 & 498 & 116 & 43 & 0.847 & 0.726 & 0.794 & 0.907 \\
Complementary & 627 & 540 & 67 & 10 & 0.902 & 0.742 & 0.861 & 0.947 \\
\textbf{Mahony} & \textbf{627} & \textbf{550} & \textbf{67} & \textbf{8} & \textbf{0.921} & \textbf{0.774} & \textbf{0.877} & \textbf{0.970} \\
Madgwick & 627 & 511 & 91 & 92 & 0.803 & 0.660 & 0.815 & 0.792 \\
Valenti & 627 & 531 & 79 & 12 & 0.898 & 0.739 & 0.847 & 0.957 \\
ESKF & 627 & 524 & 86 & 43 & 0.860 & 0.715 & 0.836 & 0.885 \\
\bottomrule
\end{tabular}

\end{table}

\begin{figure}[H]
\centering
\includegraphics[width=0.86\linewidth]{reconstruct_az/before_after_avert.png}
\caption{Raw fixed-gravity $\avert$ vs.\ reconstructed vertical acceleration
on a rotating ride. The raw trace carries a tilt-induced bias in the calm
mid-ride section that reconstruction removes.}
\label{fig:recon-before-after}
\end{figure}

\begin{figure}[H]
\centering
\includegraphics[width=0.49\linewidth]{reconstruct_az/pred_mae_bar.png}
\includegraphics[width=0.49\linewidth]{reconstruct_az/seg_f1_bar.png}
\caption{Left: prediction MAE by reconstruction. Right: segmentation F1 /
IoU-F1 by reconstruction.}
\label{fig:recon-bars}
\end{figure}

\begin{figure}[H]
\centering
\includegraphics[width=0.62\linewidth]{reconstruct_az/diagnostics_rotation.png}
\caption{Per-ride rotation energy (mean $\lvert\omega\rvert$) vs.\ the
trapezoid $\lvert$error$\rvert$ reduction from reconstruction --- the benefit
concentrates on the most-rotated rides.}
\label{fig:recon-diag}
\end{figure}

\paragraph{Metrics that help.}
Beyond MAE / F1, three signals gauge reconstruction quality: the
\emph{accel-trust} weight $w$ (fraction of the ride where the accelerometer
is a usable gravity reference), per-ride \emph{rotation energy}
$\langle\lvert\omega\rvert\rangle$ (how much reconstruction can possibly
help), and agreement of $q(t)$ with the phone's own fused \texttt{ORI}
quaternion (a free on-device reference).

\end{document}
